\clearpage
\section{Appendix: Scriptural Date and Location}
\label{sec:date-location}
% Generated. Do not edit.
% Deterministic projection of research/chronology.toml.
% schema: 5
% document: liturgy/roman-rite/1962/propers/temporal/f01-christ-the-king
% generated-by: tools/tpt proper-chronology annotations
\newcommand{\chronologyannotationclaim}[6]{#6}
\newcommand{\chronologyannotationreach}[3]{}
\newcommand{\chronologyannotationgroup}[3]{#3}
\newcommand{\chronologyannotation}[1]{%
  \ifcsname triptychchronologyannotation@#1\endcsname
    \csname triptychchronologyannotation@#1\endcsname
  \else
    \PackageError{triptych}{No chronology annotation for #1}{}%
  \fi}
% introit: Introit
\expandafter\def\csname triptychchronologyannotation@introit\endcsname{%
  \chronologyannotationgroup{historical-setting}{research-pending}{\textbf{Historical setting} -- Occasion date unresolved.}%
  \space \chronologyannotationgroup{composition}{nonuniform}{\textbf{Composition} -- No single date applies across every cited locus.}%
}
% epistle: Epistle
\expandafter\def\csname triptychchronologyannotation@epistle\endcsname{%
  \chronologyannotationgroup{composition}{disputed}{\textbf{Composition} -- disputed: \chronologyannotationreach{Col.1.12}{Col}{true}\chronologyannotationreach{Col.1.13}{Col}{true}\chronologyannotationreach{Col.1.14}{Col}{true}\chronologyannotationreach{Col.1.15}{Col}{true}\chronologyannotationreach{Col.1.16}{Col}{true}\chronologyannotationreach{Col.1.17}{Col}{true}\chronologyannotationreach{Col.1.18}{Col}{true}\chronologyannotationreach{Col.1.19}{Col}{true}\chronologyannotationreach{Col.1.20}{Col}{true}\chronologyannotationclaim{composition.epistle-to-the-colossians}{composition}{catholic-traditional-v1}{disputed}{(Philemon; Colossians; Ephesians; Philippians), 61}{A.D. 61}; \chronologyannotationreach{Col.1.12}{Col}{true}\chronologyannotationreach{Col.1.13}{Col}{true}\chronologyannotationreach{Col.1.14}{Col}{true}\chronologyannotationreach{Col.1.15}{Col}{true}\chronologyannotationreach{Col.1.16}{Col}{true}\chronologyannotationreach{Col.1.17}{Col}{true}\chronologyannotationreach{Col.1.18}{Col}{true}\chronologyannotationreach{Col.1.19}{Col}{true}\chronologyannotationreach{Col.1.20}{Col}{true}\chronologyannotationclaim{composition.epistle-to-the-colossians}{composition}{catholic-traditional-v1}{disputed}{some time between A.D. 61 and 63}{A.D. 61--63}.}%
}
% gradual: Gradual
\expandafter\def\csname triptychchronologyannotation@gradual\endcsname{%
  \chronologyannotationgroup{traditional-attribution}{disputed}{\textbf{Traditional attribution} -- disputed: \chronologyannotationreach{Ps.71.11}{Ps.2.8 Ps.28.10-11 Ps.71.1 Ps.71.8 Ps.71.11}{false}\chronologyannotationreach{Ps.71.8}{Ps.2.8 Ps.28.10-11 Ps.71.1 Ps.71.8 Ps.71.11}{false}\chronologyannotationclaim{israel.monarchy.david-traditional-era}{traditional-attribution}{catholic-traditional-v1}{disputed}{reigned from 1055 to 1015 B.C.}{David (reign in the usual chronology), B.C. 1055--1015}.}%
  \space \chronologyannotationgroup{historical-setting}{research-pending}{\textbf{Historical setting} -- Occasion date unresolved.}%
  \space \chronologyannotationgroup{composition}{preferred}{\textbf{Composition}: \chronologyannotationreach{Ps.71.11}{Ps}{true}\chronologyannotationreach{Ps.71.8}{Ps}{true}\chronologyannotationclaim{critical.psalms.latest-composition-boundary}{composition}{catholic-critical-v1}{preferred}{before the Maccabean period, around 165 B.C.; no individual psalm can be dated securely}{Before c. 165 B.C.}}%
}
% alleluia: Alleluia
\expandafter\def\csname triptychchronologyannotation@alleluia\endcsname{%
  \chronologyannotationgroup{historical-setting}{preferred}{\textbf{Historical setting}: \chronologyannotationreach{Dan.7.14}{Dan}{true}\chronologyannotationclaim{israel.exile.seventy-years}{historical-setting}{catholic-traditional-v1}{preferred}{seventy years}{The seventy years of servitude to Babylon, ``seventy years'' (duration)}.}%
  \space \chronologyannotationgroup{composition}{preferred}{\textbf{Composition}: \chronologyannotationreach{Dan.7.14}{Dan}{true}\chronologyannotationclaim{composition.book-of-daniel}{composition}{catholic-traditional-v1}{preferred}{admits 570-536 B.C. as its date of composition}{B.C. 570--536}.}%
}
% gospel: Gospel
\expandafter\def\csname triptychchronologyannotation@gospel\endcsname{%
  \chronologyannotationgroup{narrated-event}{research-pending}{\textbf{Event} -- No narrated-event date in the chronology corpus.}%
  \space \chronologyannotationgroup{composition}{disputed}{\textbf{Composition} -- disputed: \chronologyannotationreach{John.18.33}{John}{true}\chronologyannotationreach{John.18.34}{John}{true}\chronologyannotationreach{John.18.35}{John}{true}\chronologyannotationreach{John.18.36}{John}{true}\chronologyannotationreach{John.18.37}{John}{true}\chronologyannotationclaim{composition.gospel-of-john}{composition}{catholic-traditional-v1}{disputed}{from the year 90 to 100 (approximately)}{c. A.D. 90--100}; \chronologyannotationreach{John.18.33}{John}{true}\chronologyannotationreach{John.18.34}{John}{true}\chronologyannotationreach{John.18.35}{John}{true}\chronologyannotationreach{John.18.36}{John}{true}\chronologyannotationreach{John.18.37}{John}{true}\chronologyannotationclaim{composition.gospel-of-john}{composition}{catholic-traditional-v1}{disputed}{96 or one of the succeeding years}{A.D. 96--100}.}%
}
% offertory: Offertory
\expandafter\def\csname triptychchronologyannotation@offertory\endcsname{%
  \chronologyannotationgroup{traditional-attribution}{disputed}{\textbf{Traditional attribution} -- disputed: \chronologyannotationreach{Ps.2.8}{Ps.2.8 Ps.28.10-11 Ps.71.1 Ps.71.8 Ps.71.11}{false}\chronologyannotationclaim{israel.monarchy.david-traditional-era}{traditional-attribution}{catholic-traditional-v1}{disputed}{reigned from 1055 to 1015 B.C.}{David (reign in the usual chronology), B.C. 1055--1015}.}%
  \space \chronologyannotationgroup{historical-setting}{research-pending}{\textbf{Historical setting} -- Occasion date unresolved.}%
  \space \chronologyannotationgroup{composition}{preferred}{\textbf{Composition}: \chronologyannotationreach{Ps.2.8}{Ps}{true}\chronologyannotationclaim{critical.psalms.latest-composition-boundary}{composition}{catholic-critical-v1}{preferred}{before the Maccabean period, around 165 B.C.; no individual psalm can be dated securely}{Before c. 165 B.C.}}%
}
% communion: Communion
\expandafter\def\csname triptychchronologyannotation@communion\endcsname{%
  \chronologyannotationgroup{traditional-attribution}{disputed}{\textbf{Traditional attribution} -- disputed: \chronologyannotationreach{Ps.28.10}{Ps.2.8 Ps.28.10-11 Ps.71.1 Ps.71.8 Ps.71.11}{false}\chronologyannotationreach{Ps.28.11}{Ps.2.8 Ps.28.10-11 Ps.71.1 Ps.71.8 Ps.71.11}{false}\chronologyannotationclaim{israel.monarchy.david-traditional-era}{traditional-attribution}{catholic-traditional-v1}{disputed}{reigned from 1055 to 1015 B.C.}{David (reign in the usual chronology), B.C. 1055--1015}.}%
  \space \chronologyannotationgroup{historical-setting}{research-pending}{\textbf{Historical setting} -- Occasion date unresolved.}%
  \space \chronologyannotationgroup{composition}{preferred}{\textbf{Composition}: \chronologyannotationreach{Ps.28.10}{Ps}{true}\chronologyannotationreach{Ps.28.11}{Ps}{true}\chronologyannotationclaim{critical.psalms.latest-composition-boundary}{composition}{catholic-critical-v1}{preferred}{before the Maccabean period, around 165 B.C.; no individual psalm can be dated securely}{Before c. 165 B.C.}}%
}
% presentation: orientation-v2
\newcommand{\chronologyorientationrow}[2]{%
  \ifcsname triptychchronologyorientation@#1\endcsname
    \csname triptychchronologyorientation@#1\endcsname{#2}%
  \else
    \PackageError{triptych}{No chronology orientation row for #1}{}%
  \fi}
\expandafter\def\csname triptychchronologyorientation@row-f0873f1781ab3e51\endcsname#1{%
Introit: Ps 71:1 (Vulgate numbering; Hebrew 72); Gradual: Psalms 71:8,11 (Vulgate numbering; Hebrew 72); Offertory: Psalms 2:8 (Vulgate and Hebrew numbering agree); Communion: Psalm 28:10--11 (Vulgate numbering; Hebrew 29) & Traditional attribution: David's reign. #1 & Disputed: \chronologyannotationreach{Ps.2.8}{Ps.2.8 Ps.28.10-11 Ps.71.1 Ps.71.8 Ps.71.11}{false}\chronologyannotationreach{Ps.28.10}{Ps.2.8 Ps.28.10-11 Ps.71.1 Ps.71.8 Ps.71.11}{false}\chronologyannotationreach{Ps.28.11}{Ps.2.8 Ps.28.10-11 Ps.71.1 Ps.71.8 Ps.71.11}{false}\chronologyannotationreach{Ps.71.1}{Ps.2.8 Ps.28.10-11 Ps.71.1 Ps.71.8 Ps.71.11}{false}\chronologyannotationreach{Ps.71.11}{Ps.2.8 Ps.28.10-11 Ps.71.1 Ps.71.8 Ps.71.11}{false}\chronologyannotationreach{Ps.71.8}{Ps.2.8 Ps.28.10-11 Ps.71.1 Ps.71.8 Ps.71.11}{false}\chronologyannotationclaim{israel.monarchy.david-traditional-era}{traditional-attribution}{catholic-traditional-v1}{disputed}{reigned from 1055 to 1015 B.C.}{B.C. 1055--1015}\\[.25em]%
}
\expandafter\def\csname triptychchronologyorientation@row-d15f27850eb6b12a\endcsname#1{%
Introit: Ps 71:1 (Vulgate numbering; Hebrew 72); Gradual: Psalms 71:8,11 (Vulgate numbering; Hebrew 72); Offertory: Psalms 2:8 (Vulgate and Hebrew numbering agree); Communion: Psalm 28:10--11 (Vulgate numbering; Hebrew 29) & Historical setting. #1 & Occasion date unresolved\\[.25em]%
}
\expandafter\def\csname triptychchronologyorientation@row-905605cc83c0c383\endcsname#1{%
Introit: Ps 71:1 (Vulgate numbering; Hebrew 72); Gradual: Psalms 71:8,11 (Vulgate numbering; Hebrew 72); Offertory: Psalms 2:8 (Vulgate and Hebrew numbering agree); Communion: Psalm 28:10--11 (Vulgate numbering; Hebrew 29) & Composition: Psalms. #1 & \chronologyannotationreach{Ps.2.8}{Ps}{true}\chronologyannotationreach{Ps.28.10}{Ps}{true}\chronologyannotationreach{Ps.28.11}{Ps}{true}\chronologyannotationreach{Ps.71.1}{Ps}{true}\chronologyannotationreach{Ps.71.11}{Ps}{true}\chronologyannotationreach{Ps.71.8}{Ps}{true}\chronologyannotationclaim{critical.psalms.latest-composition-boundary}{composition}{catholic-critical-v1}{preferred}{before the Maccabean period, around 165 B.C.; no individual psalm can be dated securely}{Before c. 165 B.C.}\\[.25em]%
}
\expandafter\def\csname triptychchronologyorientation@row-94f74a8fe6bde9b6\endcsname#1{%
Alleluia: Daniel 7:14 & Historical setting: The seventy years of servitude to Babylon. #1 & \chronologyannotationreach{Dan.7.14}{Dan}{true}\chronologyannotationclaim{israel.exile.seventy-years}{historical-setting}{catholic-traditional-v1}{preferred}{seventy years}{``seventy years'' (duration)}\\[.25em]%
}
\expandafter\def\csname triptychchronologyorientation@row-98256791fb1a09b1\endcsname#1{%
Alleluia: Daniel 7:14 & Composition: The Book of Daniel. #1 & \chronologyannotationreach{Dan.7.14}{Dan}{true}\chronologyannotationclaim{composition.book-of-daniel}{composition}{catholic-traditional-v1}{preferred}{admits 570-536 B.C. as its date of composition}{B.C. 570--536}\\[.25em]%
}
\expandafter\def\csname triptychchronologyorientation@row-bb05ffe36adb121c\endcsname#1{%
Gospel: John 18:33--37 & Event. #1 & Date unresolved\\[.25em]%
}
\expandafter\def\csname triptychchronologyorientation@row-b5624d4c980f37b6\endcsname#1{%
Gospel: John 18:33--37 & Composition: The Gospel of St John. #1 & Disputed: \chronologyannotationreach{John.18.33}{John}{true}\chronologyannotationreach{John.18.34}{John}{true}\chronologyannotationreach{John.18.35}{John}{true}\chronologyannotationreach{John.18.36}{John}{true}\chronologyannotationreach{John.18.37}{John}{true}\chronologyannotationclaim{composition.gospel-of-john}{composition}{catholic-traditional-v1}{disputed}{from the year 90 to 100 (approximately)}{c. A.D. 90--100}; \chronologyannotationreach{John.18.33}{John}{true}\chronologyannotationreach{John.18.34}{John}{true}\chronologyannotationreach{John.18.35}{John}{true}\chronologyannotationreach{John.18.36}{John}{true}\chronologyannotationreach{John.18.37}{John}{true}\chronologyannotationclaim{composition.gospel-of-john}{composition}{catholic-traditional-v1}{disputed}{96 or one of the succeeding years}{A.D. 96--100}\\[.25em]%
}
\expandafter\def\csname triptychchronologyorientation@row-e679cfff41e9473d\endcsname#1{%
Epistle: Colossians 1:12--20 & Composition: The Epistle to the Colossians. #1 & Disputed: \chronologyannotationreach{Col.1.12}{Col}{true}\chronologyannotationreach{Col.1.13}{Col}{true}\chronologyannotationreach{Col.1.14}{Col}{true}\chronologyannotationreach{Col.1.15}{Col}{true}\chronologyannotationreach{Col.1.16}{Col}{true}\chronologyannotationreach{Col.1.17}{Col}{true}\chronologyannotationreach{Col.1.18}{Col}{true}\chronologyannotationreach{Col.1.19}{Col}{true}\chronologyannotationreach{Col.1.20}{Col}{true}\chronologyannotationclaim{composition.epistle-to-the-colossians}{composition}{catholic-traditional-v1}{disputed}{(Philemon; Colossians; Ephesians; Philippians), 61}{A.D. 61}; \chronologyannotationreach{Col.1.12}{Col}{true}\chronologyannotationreach{Col.1.13}{Col}{true}\chronologyannotationreach{Col.1.14}{Col}{true}\chronologyannotationreach{Col.1.15}{Col}{true}\chronologyannotationreach{Col.1.16}{Col}{true}\chronologyannotationreach{Col.1.17}{Col}{true}\chronologyannotationreach{Col.1.18}{Col}{true}\chronologyannotationreach{Col.1.19}{Col}{true}\chronologyannotationreach{Col.1.20}{Col}{true}\chronologyannotationclaim{composition.epistle-to-the-colossians}{composition}{catholic-traditional-v1}{disputed}{some time between A.D. 61 and 63}{A.D. 61--63}\\[.25em]%
}
\expandafter\def\csname triptychchronologyorientation@row-1dc5d7a61ec35b95\endcsname#1{%
Introit: Apoc 1:6, Apoc 5:12 & Composition: The Apocalypse of St John. #1 & Disputed: \chronologyannotationreach{Apoc.1.6}{Apoc}{true}\chronologyannotationreach{Apoc.5.12}{Apoc}{true}\chronologyannotationclaim{composition.apocalypse}{composition}{catholic-traditional-v1}{disputed}{the reign of Domitian ... A.D. 81-96}{A.D. 81--96}; \chronologyannotationreach{Apoc.1.6}{Apoc}{true}\chronologyannotationreach{Apoc.5.12}{Apoc}{true}\chronologyannotationclaim{composition.apocalypse}{composition}{catholic-traditional-v1}{disputed}{from the year 90 to 100 (approximately)}{c. A.D. 90--100}\\[.25em]%
}

The table distinguishes the date of an attributed person, a historical setting, a narrated event and the writing of a book. A date for one subject does not date the others. Its passage references identify the complete verse envelopes used for chronology; the actual chant selections remain those printed above. The three Psalm rows share passages because attribution, occasion and composition answer different questions. The numbered discussions following the table identify the historical evidence and its limits.

\begin{chronologyorientationtable}
\chronologyorientationrow{row-ef4b73ec52fbafba}{Israel and Judah in the received Davidic frame; see H1.}
\chronologyorientationrow{row-6c43dd5e0531017f}{Royal prayer, Sion and the nations, and the storm world of the Psalter; precise occasions unknown. See H1.}
\chronologyorientationrow{row-ff280e363c1a762b}{Israel's Psalter; individual writing places uncertain. See H1.}
\chronologyorientationrow{row-94f74a8fe6bde9b6}{Daniel's Babylonian royal setting; the court is a vision. See H2.}
\chronologyorientationrow{row-98256791fb1a09b1}{Babylon in the received authorial setting; formation disputed. See H2.}
\chronologyorientationrow{row-bb05ffe36adb121c}{Jesus before Pilate in the governor's hall at Jerusalem, John 18:28,33. See H3.}
\chronologyorientationrow{row-b5624d4c980f37b6}{Ephesus traditionally; writing place and formation disputed. See H3.}
\chronologyorientationrow{row-e679cfff41e9473d}{Rome traditionally; authorship and prison location disputed. Colossae is the recipients' city. See H4.}
\chronologyorientationrow{row-1dc5d7a61ec35b95}{Patmos in the received seer's setting; heaven is the scene of the vision. See H5.}
\end{chronologyorientationtable}

\subsection{H1. The Psalms: attribution, occasion and writing}
\label{history-psalms}
The Davidic reference era comes from John Corbett's \emph{King David} in the \emph{Catholic Encyclopedia}, IV, opening paragraph. Corbett calls it the usual chronology and immediately records proposals that place the reign thirty to fifty years later. The displayed range dates the attributed figure. It neither locates every poem within his reign nor supplies the date of a particular performance.

The three Psalms reach David through different evidence. Psalm 28's title names him. Psalm 2 has no such name in the inspected title, but Acts 4:24--28 gives its Davidic attribution in apostolic reception. Psalm 71's title concerns Solomon; Bellarmine understands its opening as David's prayer for the newly anointed king. That is his interpretation of the speaker, not the wording of a title that explicitly names David as author. Augustine begins directly from Christ, the true Solomon. The difference belongs to the reception of the psalm and does not establish a new composition date.

The historical occasions remain uncertain. Psalm 28's tabernacle-completion title does not identify a securely dated building project; its Libanus and Cades belong to the psalm's storm world. Psalm 2 names Sion and the nations without dating the appointment of its royal speaker. Bellarmine's accession setting for Psalm 71 is an exegetical identification. No shared precise place of writing follows from these different settings.

The composition boundary comes from the USCCB NABRE introduction to Psalms, in its discussion of superscriptions and dating. It describes a pre-Maccabean limit alongside the inability to date individual psalms securely. Thus ``before'' remains part of the answer. The boundary is not an estimate that these poems were written near its endpoint. The complete contemporary introduction was independently checked in the reviewed research; its historical judgment is distinct from Corbett's conventional chronology and the received attributions.\footnote{Corbett, \emph{King David}, CE IV, opening; USCCB, \emph{Psalms: Introduction}, superscriptions and dating discussion; Bellarmine, Ps 71, opening; Augustine, Ps 71.1--2.}

\subsection{H2. Daniel: Babylon and the heavenly court}
\label{history-daniel}
Daniel 7:1 places the vision in the first year of Baltasar, king of Babylon. The heavenly court is the visionary scene, not a terrestrial place of composition. The table's seventy years describes Babylonian servitude in the received historical frame, supported by Jeremiah 25:11--12. A duration cannot be used as a calendar year for the Son of Man's reception of dominion.

Francis Gigot's \emph{Book of Daniel}, CE IV, authorship discussion, reports the displayed traditional composition range and the competing theory locating the book's formation in the Maccabean crisis under Antiochus IV. Those are different accounts of the book's history. Jerome's rejection of a fulfillment exhausted by the Maccabees concerns the meaning and endurance of the promised kingdom; it is not a modern critical judgment about the date of composition. The interpretation of the prophecy and the circumstances of the book's formation remain distinguishable questions.\footnote{Gigot, \emph{Book of Daniel}, authorship discussion; Dan 7:1,13--14,17--18,27; Jer 25:11--12; Jerome, \emph{Commentary on Daniel}, VII, on 13--18.}

\subsection{H3. John: the trial and the Gospel's formation}
\label{history-john}
The narrated scene is at Jerusalem, in the governor's hall, in the morning of the Passion sequence. John 18:28 and 33 locate the interrogation; John 19:14--18 carries the narrative toward sentence and execution. A chronological binding for the Crucifixion cannot simply be attached to this distinct scene. The table therefore retains an unresolved date for the interrogation. Howlett's \emph{Biblical Chronology}, CE III, Passion discussion, preserves several alternative years and expressly declines certainty; its disagreement supplies context rather than a selected date for this row.

The received attribution names John son of Zebedee, and Ephesus is the traditional writing place. The \emph{Catholic Encyclopedia} articles \emph{New Testament}, XIV, dating discussion, and \emph{Gospel of St John}, VIII, circumstances of composition, supply the two displayed late composition ranges. The latter acknowledges the lack of certain historical information and proposes a year within the last decade rather than an exact completion date established by contemporary documentation.

The USCCB NABRE introduction to John distinguishes the eyewitness testimony underlying the tradition from authorship of the finished Gospel. It reports that most modern scholars do not confidently identify that witness with Zebedee's son. Its account of successive editing points to chapter 21, the adapted prologue hymn and discontinuities in the farewell discourse. Neither the number nor identities of all contributors are established. Alongside Ephesus it reports proposals for Syria, possibly Antioch, and Alexandria. These concern the writing and formation of the Gospel; the trial remains at Jerusalem. The displayed composition ranges are not reassigned as exact dates for particular editorial layers.\footnote{CE, \emph{Gospel of St John}, circumstances of composition; \emph{New Testament}, dating discussion; Howlett, \emph{Biblical Chronology}, Passion discussion. USCCB, \emph{John: Introduction}, paragraphs beginning ``Critical analysis,'' ``Other difficulties,'' ``The gospel contains,'' and ``The final editing,'' inspected 20 October 2026.}

\subsection{H4. Colossians: the prison and the recipients}
\label{history-colossians}
Pauline authorship and Roman captivity form the received setting used by the displayed chronology. The \emph{Catholic Encyclopedia}'s \emph{St Paul}, XI, chronology list, supplies the single-year answer; \emph{Epistle to the Colossians}, IV, opening and captivity discussion, supplies the wider range. Different precision need not imply a direct contradiction. The figures date the letter's writing, not the creation or reconciliation it proclaims. Colossae is the recipients' city and cannot be substituted for the prison's location.

The USCCB NABRE introduction notes that the letter does not identify the prison and that Paul's several imprisonments complicate dating. It reports Roman, Caesarean and Ephesian proposals, and also the possibility of a pupil or follower writing in Paul's name. It does not decide between that account and direct Pauline authorship or supply the follower with a definite city and date. The traditional Roman chronology therefore remains identified as such, beside the distinct uncertainties of authorship and imprisonment.\footnote{CE, \emph{St Paul}, chronology list; \emph{Epistle to the Colossians}, opening and captivity discussion; USCCB, \emph{Colossians: Introduction}, paragraphs 2 and 5, especially ``Paul wrote,'' inspected 20 October 2026.}

\subsection{H5. Apocalypse: the seer's place and the vision}
\label{history-apocalypse}
The two Apocalypse passages in the Introit share a book-composition answer. The Psalm verse does not share that answer. The \emph{Catholic Encyclopedia}'s \emph{Apocalypse}, I, time-and-place section, presents Patmos and testimony connecting the work with the end of Domitian's reign. The whole reign provides a broad bound, not a statement that each year is equally probable. The \emph{New Testament} article supplies the alternative approximate late range and acknowledges an earlier Neronian proposal. The selected late chronology does not exhaust the history of scholarly positions.

Patmos belongs to the received seer's setting; the heavenly throne room belongs to the vision. Neither should be mistaken for the other. Apocalypse 1 situates the witness amid suffering and patient endurance, while chapter 5 displays worship around the Lamb. That literary relation supports the study's interpretation without assigning a historical date to the eternal kingship being acclaimed.\footnote{CE, \emph{Apocalypse}, time-and-place section; \emph{New Testament}, dating discussion; Apoc 1:1--20; 5:1--14.}
